\section{Closing}
\label{sec:closing}

\subsection{The chain, in one page}

The mathematics begins with an expectation over trajectories
\eqref{eq:objective} and makes its gradient estimable without a model of the
environment \eqref{eq:reinforce}. The policy gradient theorem then replaces the
trajectory return with a state--action value, at the price of weighting by the
policy's own state distribution \eqref{eq:pgt}. Substituting a learned value
gives the actor--critic family, and the choice of how to train that value decides
the stability of everything above it: bootstrapping buys sample efficiency and
risks divergence.

In the deterministic branch, the same identity survives without the action
integral \eqref{eq:dpg}, which makes continuous control tractable; that branch is
where the deep deterministic policy gradient lives, and it is not the branch that
produced AlphaGo.

The AlphaGo branch replaces the gradient estimator entirely. Policy improvement
becomes tree search; policy evaluation becomes regression on the outcome of games
played by the improved policy; the network is refitted to both. The result is an
approximate policy iteration whose two operators are a learned function and a
search built on top of it, and its theoretical frame is Expert Iteration. The
bandit results supply the selection rule and the only real guarantees in the
system, and the guarantees stop where the learned prior begins.

The engineering has two hard problems, and they are the same problem seen twice:
a tree that owns its children in a language where recursive ownership is
expensive to express, and a single accelerator that must serve many cheap,
irregular searches. The answers are a flat arena with integer indices, and many
independent trees feeding one batching service. Both answers trade a
compile-time guarantee for a stated invariant, and both are worth the trade
precisely because the invariants are small enough to write down
(Section~\ref{sec:invariants}).

\subsection{What is guaranteed, and what is measured}

The system that this monograph describes is an empirically justified system. It
contains, at its centre, three components with no convergence proof: PUCT with
learned priors, approximate policy iteration with deep function approximation,
and the self-play loop itself (Table~\ref{tab:proven}). Every one of them worked,
in the primary sources and in several independent replications, and none of them
comes with a theorem.

That is not a defect to be fixed; it is the normal state of deep reinforcement
learning, and it has one practical consequence that shapes the engineering more
than any theoretical consideration: correctness must be established against
known facts rather than by appeal to theory. Gomoku is unusually generous in this
respect. Positions with exact values and forced-win sequences with mechanical
verification are available without any learning, and they can be used to test
the search, the prover's soundness, and the sign conventions before a single
gradient is computed. A system that passes those tests is still not proven
correct, but it is anchored to ground truth at the points where defects are most
likely and most expensive.

\subsection{The decisions that are expensive to reverse}

Three choices in this monograph are worth more thought than the rest, because
they are the ones that are painful to change later.

\emph{The tree representation.} A flat arena with integer indices is not
elegant, and it gives up a compile-time guarantee. It is chosen because the
alternative representations either do not compile or pay runtime costs in the
hottest loop in the system. Everything that touches the tree -- traversal,
backup, reuse across moves, serialization for debugging -- is written against
this representation, so changing it later is a rewrite of the search crate.

\emph{Ownership of the model.} Making the network a single owner on a single
thread, with plain messages as the interface, removes an entire class of
concurrency problems from the design rather than managing them. The alternative
-- shared weights behind a lock -- looks simpler in the first ten lines of code
and is worse in every line after that.

\emph{What the replay window stores.} Storing games rather than encoded positions
keeps the encoding decisions mutable and reduces memory by orders of magnitude.
It also means the corpus remains valid while the plane layout is still being
experimented with, which is exactly the period during which the corpus is
growing.

\subsection{The order of the upgrades}

Section~\ref{sec:proven} argues that the guarantees live in the skeleton, and
Section~\ref{sec:scaling} transfers the published constants to a smaller problem.
The natural next step is to change things, and the order matters, for one
structural reason: some upgrades change the distribution of the training data and
others do not. A change of the data distribution cannot be evaluated against a
baseline whose data came from a different distribution, so such changes must come
after a clean, working baseline exists, and must be made one at a time.

The upgrades that leave the data distribution alone can be made early: growing
the network, adding pooling to the heads, moving from phased to continuous
execution, and improving the evaluation setups. The upgrades that change the data
distribution -- sampling some moves with a smaller search budget, adding
auxiliary prediction targets, changing the temperature schedule, changing the
simulation count -- must wait, and must be measured against a frozen baseline by
playing games rather than by watching the loss.

The general rule is the one every part of this monograph has followed: decide
what can be decided by argument, measure what cannot, and keep the set of things
being changed at once to one.

\subsection{Reading the primary sources}

The three papers reward being read in sequence rather than separately, because
the deletions of Table~\ref{tab:removals} are the argument. For the
mathematics, the policy gradient theorem and the deterministic policy gradient
papers are short and worth their length; Expert Iteration names the two moving
parts of the loop; and the UCT and PUCT papers are the two ends of the search
argument, including the place where its guarantee stops.
